\uuid{8I0p}
\titre{Limite en l'infini}
\auteur{bodin ; hanani ; mzari}
\organisation{exo7}
\licence{CC BY-NC-SA 4.0}
\datecreate{2020-01-20}
\niveau{L1}
\module{Analyse}
\chapitre{Continuité, limite et étude de fonctions réelles}
\sousChapitre{Limite de fonctions}
\difficulte{2}
\theme{Limites des fonctions réelles, limite en l'infini}
\exo7qcmid{lille-1.limites.012}

\contenu{
\begin{question}
\qclassification{123.03}
\qtags{motcle=limite en l'infini}
Soit  $f$ une fonction définie sur $\Rr$. Quelles sont les assertions vraies ?
\begin{answers}

    \bad{$\lim_{x\to +\infty} f(x)=l \, (l\in \Rr)$  si et seulement si $ \forall \epsilon >0,  \exists A >0, \forall x \in \Rr, |f(x)-l|<\epsilon \Rightarrow x>A$}

    \good{$\lim_{x\to +\infty} f(x)=l \, (l\in \Rr)$  si et seulement si $\forall \epsilon >0,  \exists A >0, \forall x \in \Rr, x\ge A \Rightarrow |f(x)-l|\le \epsilon$}

    \good{$\lim_{x\to -\infty} f(x)=+\infty$  si et seulement si $ \forall A >0,  \exists B <0, \forall x \in \Rr, x\le B \Rightarrow f(x)\ge A$}

    \bad{$\lim_{x\to -\infty} f(x)=-\infty$  si et seulement si $  \exists B<0, \forall A < 0,   \forall x \in \Rr, x < B \Rightarrow f(x) < A$}

\end{answers}
\begin{explanations}
Voir la définition d'une limite en $+\infty$ ou $-\infty$ vers une valeur finie ou infinie :

$\lim_{x\to +\infty} f(x)=l \, (l\in \Rr)$  si et seulement si $\forall \epsilon >0,  \exists A >0, \forall x \in \Rr, x\ge A \Rightarrow |f(x)-l|\le \epsilon$

$\lim_{x\to -\infty} f(x)=+\infty$  si et seulement si $ \forall A >0,  \exists B <0, \forall x \in \Rr, x\le B \Rightarrow f(x)\ge A$
\end{explanations}
\end{question}
}
